\chapter{LANGUAGE DESIGN AND GRAMMAR}

\section{Types}
The language provides four primitive types, chosen to cover the core domains
of numeric computation, logical control flow, and text manipulation while
keeping the type system tractable and the semantics well-defined.

\subsection{Integer}
The Integer type represents whole numbers such as \texttt{0}, \texttt{7}, and
\texttt{42}. Integer literals are written as one or more decimal digits without
quotation marks or a decimal point.

\subsection{Float}
The Float type represents real-valued numbers written with a decimal point, such
as \texttt{3.14} or \texttt{0.5}. Float values are used exclusively with the
dot-suffixed arithmetic operators (\texttt{+.}, \texttt{-.}, \texttt{*.},
\texttt{/.}), following the Caml~Light convention. Integer division (\texttt{/})
performs floor division and returns \texttt{Integer}: for example,
\texttt{10 / 3} evaluates to \texttt{3 : Integer}.

The language intentionally does not perform implicit numeric promotion between
\texttt{Integer} and \texttt{Float} values. This choice follows directly from
the design requirement of ``no overloading with the usual precedence of
multiplication and division over addition and subtraction.'' In a language
where a single \texttt{+} operator accepts both
$(\textit{Integer}, \textit{Integer}) \rightarrow \textit{Integer}$ and
$(\textit{Integer}, \textit{Float}) \rightarrow \textit{Float}$, the operator
effectively carries multiple type signatures, since the compiler must inspect
operand types to determine the result type. This constitutes operator
overloading, even when described as ``numeric promotion'' or ``widening
coercion.'' To eliminate that ambiguity, the language follows the Caml~Light
convention of using separate operator tokens for integer and floating-point
arithmetic: \texttt{+}, \texttt{-}, \texttt{*}, and \texttt{/} operate only on
\texttt{Integer} operands, whereas \texttt{+.}, \texttt{-.}, \texttt{*.}, and
\texttt{/.} operate only on \texttt{Float} operands. Mixed-type arithmetic such
as \texttt{1 + 2.5} is therefore rejected by the type checker. Under this
design, each operator token uniquely determines both operand and result types,
so operator meaning never depends on post hoc type resolution.

\subsection{Boolean}
The Boolean type represents logical truth values. The two Boolean literals are
\texttt{true} and \texttt{false}. Boolean values are used in conditions and as
results of comparison expressions.

\subsection{String}
The String type represents sequences of characters enclosed in double quotes,
for example \texttt{"hello"}. Strings can be assigned to variables and printed
through the built-in \texttt{print()} function. The surrounding quotes are
stripped by the parser when the \texttt{Literal} node is created, so the value
stored internally is the bare character sequence without surrounding quotes. At
runtime, the built-in \texttt{print()} function re-adds double quotes around
string values, so \texttt{print("plc")} outputs \texttt{"plc" : String}.

\subsection{Void}
Although the language exposes exactly four programmer-visible types, the
implementation introduces a fifth value, \texttt{Void}, within the
\texttt{LanguageType} enumeration. Its sole purpose is to provide a
well-defined return type for functions that contain no \texttt{return}
statement, allowing the symbol table to record a consistent entry for every
function regardless of whether it produces a value. \texttt{Void} cannot appear
in source code and is never visible to the programmer; it is a bookkeeping
device used internally by the type checker. Therefore, \texttt{Void} does not
extend the programmer-visible source-level type system.

\section{Static Typing}
The language employs static typing, meaning that type errors are detected at
type-checking time rather than at runtime. No explicit variable type
declarations are required. Instead, the type checker infers types from literals,
assignments, operations, function calls, and return statements. Once a variable
type is established, later assignments must remain consistent with that type.
For example, the following program is rejected at the type-checking stage before
any execution occurs:

\begin{lstlisting}
x = 5;
x = "hello";
\end{lstlisting}

\noindent The type checker rejects the second assignment with:
\begin{lstlisting}
[line 2, col 1] TypeError: Cannot assign String to variable 'x' of type Integer
\end{lstlisting}

\noindent This contrasts with dynamically-typed languages, where the same
program would run without error until (or unless) a type-dependent operation
was attempted.

\section{Lexeme Grammar}
\label{sec:lexeme-grammar}
The \textbf{lexeme grammar} operates at the character level: it defines the
regular expressions the lexer uses to scan raw source characters into lexeme
strings. The notation below uses \texttt{+} (one or more), \texttt{*} (zero or
more), \texttt{|} (alternation), and \texttt{[\ldots]} (character class).

\begin{lstlisting}
(* Character classes *)
digit   ::=  [0-9]
letter  ::=  [a-zA-Z] | '_'

(* Literals *)
INT_LITERAL     ::=  digit+
FLOAT_LITERAL   ::=  digit+ '.' digit+
BOOL_LITERAL    ::=  'true' | 'false'
STRING_LITERAL  ::=  '"' [^"\n]* '"'

(* Identifiers and keywords *)
IDENTIFIER  ::=  letter (letter | digit)*
                 -- where the lexeme is not in the KEYWORDS set
KEYWORD     ::=  'if' | 'else' | 'while' | 'def' | 'return'

(* Arithmetic operators -- integer and float variants *)
PLUS        ::=  '+'     PLUS_DOT    ::=  '+.'
MINUS       ::=  '-'     MINUS_DOT   ::=  '-.'
STAR        ::=  '*'     STAR_DOT    ::=  '*.'
SLASH       ::=  '/'     SLASH_DOT   ::=  '/.'

(* Comparison and assignment *)
EQUAL_EQUAL ::=  '=='
BANG_EQUAL  ::=  '!='
ASSIGN      ::=  '='

(* Delimiters *)
LPAREN ::= '('   RPAREN    ::= ')'
LBRACE ::= '{'   RBRACE    ::= '}'
COMMA  ::= ','   SEMICOLON ::= ';'
\end{lstlisting}

\noindent\textbf{Maximal-munch rule.}
The lexer applies the maximal-munch (longest-match) principle: at each
position it consumes the longest character sequence that matches any rule.
This ensures, for example, that \texttt{==} is scanned as one
\texttt{EQUAL\_EQUAL} token rather than two \texttt{ASSIGN} tokens, and that
\texttt{+.} is scanned as one \texttt{PLUS\_DOT} token rather than a
\texttt{PLUS} followed by a stray dot (unless the \texttt{.} is followed by a
digit, in which case it starts a \texttt{FLOAT\_LITERAL} and the preceding
\texttt{+} is emitted as \texttt{PLUS}).

\noindent\textbf{Keyword and boolean priority.}
\texttt{IDENTIFIER} and \texttt{KEYWORD} share the same character-level
pattern. The lexer resolves ambiguity by scanning an identifier-shaped lexeme
first, then consulting the \texttt{KEYWORDS} map. If the lexeme matches a
keyword entry (including \texttt{true} and \texttt{false}, which map to
\texttt{BOOL\_LITERAL}), the keyword token type is emitted; otherwise an
\texttt{IDENTIFIER} token is produced. The name \texttt{print} is not a
keyword and is therefore always scanned as an \texttt{IDENTIFIER}.

\section{Token Grammar}
\label{sec:token-grammar}
The \textbf{token grammar} maps each lexeme class produced by the lexeme
grammar to a named token type. Table~\ref{tab:token-specification} lists all
token types recognised by the language together with their categories; the
lexeme patterns are those defined formally in
Section~\ref{sec:lexeme-grammar}.

\begin{table}[H]
\centering
\caption{Token Grammar: token types and their categories}
\label{tab:token-specification}
\begin{tabular}{lll}
\toprule
Token Name & Category & Lexeme(s) \\
\midrule
\texttt{INT\_LITERAL}   & Integer literal   & one or more decimal digits, e.g.\ \texttt{0}, \texttt{10}, \texttt{42} \\
\texttt{FLOAT\_LITERAL} & Float literal     & digits with decimal point, e.g.\ \texttt{3.14}, \texttt{20.5} \\
\texttt{BOOL\_LITERAL}  & Boolean literal   & \texttt{true}, \texttt{false} \\
\texttt{STRING\_LITERAL}& String literal    & double-quoted text, e.g.\ \texttt{"hello"} \\
\texttt{IDENTIFIER}     & Identifier        & letter or \texttt{\_} followed by letters, digits, or \texttt{\_} \\
\midrule
\texttt{IF}             & Keyword           & \texttt{if} \\
\texttt{ELSE}           & Keyword           & \texttt{else} \\
\texttt{WHILE}          & Keyword           & \texttt{while} \\
\texttt{DEF}            & Keyword           & \texttt{def} \\
\texttt{RETURN}         & Keyword           & \texttt{return} \\
\midrule
\texttt{PLUS}           & Integer arithmetic & \texttt{+} \\
\texttt{MINUS}          & Integer arithmetic & \texttt{-} \\
\texttt{STAR}           & Integer arithmetic & \texttt{*} \\
\texttt{SLASH}          & Integer arithmetic & \texttt{/} \\
\texttt{PLUS\_DOT}      & Float arithmetic  & \texttt{+.} \\
\texttt{MINUS\_DOT}     & Float arithmetic  & \texttt{-.} \\
\texttt{STAR\_DOT}      & Float arithmetic  & \texttt{*.} \\
\texttt{SLASH\_DOT}     & Float arithmetic  & \texttt{/.} \\
\texttt{EQUAL\_EQUAL}   & Comparison        & \texttt{==} \\
\texttt{BANG\_EQUAL}    & Comparison        & \texttt{!=} \\
\texttt{ASSIGN}         & Assignment        & \texttt{=} \\
\midrule
\texttt{LPAREN}         & Delimiter         & \texttt{(} \\
\texttt{RPAREN}         & Delimiter         & \texttt{)} \\
\texttt{LBRACE}         & Delimiter         & \texttt{\{} \\
\texttt{RBRACE}         & Delimiter         & \texttt{\}} \\
\texttt{COMMA}          & Delimiter         & \texttt{,} \\
\texttt{SEMICOLON}      & Delimiter         & \texttt{;} \\
\texttt{EOF}            & End of input      & (end of source) \\
\bottomrule
\end{tabular}
\end{table}

\section{Context-Free Grammar}
The language grammar can be described by the context-free grammar
$G = (V_N, V_T, P, S)$, where $V_N$ is the set of non-terminals, $V_T$ is the
set of terminals, $P$ is the set of production rules, and $S$ is the start
symbol \texttt{program}. The grammar is implemented through recursive-descent
parsing functions.
For readability, the presentation below mixes classical CFG productions with
EBNF shorthand in selected rules; both notations correspond to the same parser
implementation.

\subsection{Program Structure}
\begin{lstlisting}
program   -> stmt_list EOF
statement -> assignment
           | if_stmt
           | while_stmt
           | func_def
           | return_stmt
           | expr SEMICOLON
stmt_list -> statement stmt_list | epsilon
\end{lstlisting}

\subsection{Expressions}
An \textbf{expression} is any syntactic construct that evaluates to a value of
one of the four programmer-visible types: \texttt{Integer}, \texttt{Float},
\texttt{Boolean}, or \texttt{String}. Concretely, an expression is one of:

\begin{itemize}
  \item a literal value (\texttt{INT\_LITERAL}, \texttt{FLOAT\_LITERAL},
        \texttt{BOOL\_LITERAL}, or \texttt{STRING\_LITERAL});
  \item a variable reference (\texttt{IDENTIFIER});
  \item a function call (\texttt{IDENTIFIER} followed by a parenthesised
        argument list);
  \item a unary arithmetic negation (\texttt{MINUS} or \texttt{MINUS\_DOT}
        applied to an expression);
  \item a binary arithmetic operation (addition, subtraction, multiplication,
        or division, in integer or float variants);
  \item a comparison (\texttt{EQUAL\_EQUAL} or \texttt{BANG\_EQUAL}) between
        two arithmetic expressions, yielding a \texttt{Boolean}; or
  \item a parenthesised expression.
\end{itemize}

\noindent Expressions are the only constructs that carry a type; statements do
not produce values. The grammar rules below define the precise structure and
operator-precedence hierarchy of expressions:

\begin{lstlisting}
expr        -> comparison | additive
comparison  -> additive (EQUAL_EQUAL | BANG_EQUAL) additive
additive    -> term ((PLUS | MINUS | PLUS_DOT | MINUS_DOT) term)*
term        -> factor ((STAR | SLASH | STAR_DOT | SLASH_DOT) factor)*
factor      -> MINUS factor              (* integer unary negation *)
             | MINUS_DOT factor          (* float unary negation *)
             | INT_LITERAL
             | FLOAT_LITERAL
             | STRING_LITERAL
             | BOOL_LITERAL
             | IDENTIFIER
             | IDENTIFIER LPAREN args? RPAREN
             | LPAREN expr RPAREN
\end{lstlisting}

The \texttt{comparison} non-terminal is introduced so that the parse-tree
renderer can label equality and inequality expressions explicitly. When an
\texttt{expr} contains no \texttt{==} or \texttt{!=}, it expands directly to an
\texttt{additive}; otherwise, it expands to a \texttt{comparison} node whose
two operands are themselves \texttt{additive} sub-trees. Comparisons remain
non-associative (one comparison per expression) and produce a
\texttt{Boolean} result.

The operators \texttt{+}, \texttt{-}, \texttt{*}, and \texttt{/} operate
exclusively on \texttt{Integer} operands. The dot-suffixed operators
\texttt{+.}, \texttt{-.}, \texttt{*.}, and \texttt{/.} operate exclusively on
\texttt{Float} operands. This follows the Caml~Light convention of using
distinct operator tokens for integer and float arithmetic, eliminating the need
for implicit numeric promotion. Unary negation is similarly split:
\texttt{-x} negates an \texttt{Integer} (represented internally as the
equivalent subtraction form \texttt{0 - x}), and
\texttt{-.x} negates a \texttt{Float} (represented internally as
\texttt{0.0 -. x}).

\textbf{Integer division and division by zero.}
\texttt{/} is floor (integer) division: both operands must be \texttt{Integer},
the result is \texttt{Integer}, and the value is $\lfloor a / b \rfloor$.
For example, \texttt{7 / 2} evaluates to \texttt{3 : Integer}.
Float division uses the distinct operator \texttt{/.} and requires two
\texttt{Float} operands.
Division by zero (\texttt{x / 0} or \texttt{x /. 0.0}) is detected at runtime
and raises a \texttt{RuntimeError}.

\subsection{Boolean Expressions}
Comparison operators \texttt{==} and \texttt{!=} are the lowest-precedence
binary operators and produce \texttt{Boolean} values. As they are part
of the general \texttt{expr} rule, Boolean expressions may appear both as
control-flow conditions and as ordinary expressions. Their inferred type is
\texttt{Boolean}, allowing comparison results to be assigned to variables,
passed as arguments, or printed:

\begin{lstlisting}
flag = (x == y);
print(a != b);
\end{lstlisting}

\noindent Here \texttt{flag} is inferred as \texttt{Boolean}, and
\texttt{print(a != b)} outputs either \texttt{true : Boolean} or
\texttt{false : Boolean}.

\textbf{Strict arithmetic operands.}
The specification states that \texttt{==} and \texttt{!=} operate between two
arithmetic expressions. In this language, ``arithmetic'' is interpreted
strictly as \texttt{Integer} and \texttt{Float}, consistent with the
no-overloading rule. Both operands must also share the same numeric type, with
no implicit coercion. Therefore
\texttt{Integer == Integer} and \texttt{Float == Float} are well-typed;
\texttt{Integer == Float}, \texttt{Boolean == Boolean}, and
\texttt{String == String} are all rejected at type-check time.

\subsection{Statements}
\begin{lstlisting}
assignment  -> IDENTIFIER ASSIGN expr SEMICOLON
if_stmt     -> IF LPAREN expr RPAREN block else_part
else_part   -> ELSE block | epsilon
while_stmt  -> WHILE LPAREN expr RPAREN block
return_stmt -> RETURN expr SEMICOLON
block       -> LBRACE stmt_list RBRACE
stmt_list   -> statement stmt_list | epsilon
\end{lstlisting}

\subsubsection*{Design note: \texttt{print} as a built-in callable}
\texttt{print} is implemented as a normal function-call expression rather than
as a dedicated statement form. Lexically, \texttt{print} is tokenised as an
\texttt{IDENTIFIER}; syntactically, \texttt{print(x)} is parsed through the same
\texttt{FunctionCall} path used for user-defined calls. Semantically it is a
unary built-in of type \texttt{T $\rightarrow$ Void}: it accepts one argument
of any programmer-visible type, emits \texttt{value : Type} at runtime, and
returns the internal \texttt{Void} type. As with any void-returning call,
\texttt{print(...)} may be used as a standalone expression statement but not as
the right-hand side of an assignment.
\texttt{Void} is internal bookkeeping only: it never appears as a source-level
type annotation or source-level type literal.

\subsection{Function Definitions and Calls}
\begin{lstlisting}
func_def    -> DEF IDENTIFIER LPAREN params RPAREN block
params      -> IDENTIFIER param_tail | epsilon
param_tail  -> COMMA IDENTIFIER param_tail | epsilon
args        -> expr arg_tail | epsilon
arg_tail    -> COMMA expr arg_tail | epsilon
\end{lstlisting}

The parser implementation accepts the equivalent EBNF shorthand with optional
parameter and argument lists. Concretely, this means the grammar allows both
\texttt{def f() \{\}} and \texttt{f()}. The corresponding parse-tree renderer makes the empty case
explicit by emitting a single \texttt{epsilon} leaf as the only child of an
empty \texttt{params} or \texttt{args} node, so an empty parameter or argument
list is still labelled in the tree.

Function calls use value parameter passing, that is, call-by-value. The
argument values are evaluated before the call and copied into the function's
local environment.

The built-in \texttt{print()} function uses exactly the same surface syntax as
other calls:

\begin{lstlisting}
print(x);
print(3.14);
\end{lstlisting}

\textbf{Return type uniformity.}
A function's return type is inferred from its \texttt{return} statements.
All \texttt{return} expressions within a function body must yield the same
type; mixed return types are a static type error. A function with no
\texttt{return} statement is assigned the special internal type \texttt{Void}
and may only be invoked as a standalone statement, not used as an expression.
Assigning the result of a \texttt{Void} function to a variable (e.g.\
\texttt{x = f();} where \texttt{f} returns nothing) is rejected by the type
checker. Passing a \texttt{Void}-returning function as an argument ---
for example \texttt{print(f())} where \texttt{f} returns nothing --- is also
rejected, since \texttt{Void} is not a printable value.

\section{Operator Precedence and Associativity}
Operator precedence is encoded structurally in the recursive-descent parser.
Comparison is lowest, addition and subtraction are next, and multiplication and
division have the highest precedence among binary operators. Arithmetic
operators are left-associative.

\begin{table}[H]
\centering
\caption{Operator Precedence (1 = lowest, 3 = highest) and Associativity}
\begin{tabular}{lll}
\toprule
Precedence & Operators & Associativity \\
\midrule
1 (lowest) & \texttt{==}, \texttt{!=} & non-associative (one comparison per expression) \\
2 & \texttt{+}, \texttt{-}, \texttt{+.}, \texttt{-.} & left-associative \\
3 (highest) & \texttt{*}, \texttt{/}, \texttt{*.}, \texttt{/.} & left-associative \\
\bottomrule
\end{tabular}
\end{table}

\begin{table}[H]
\centering
\caption{Formal arithmetic and comparison operator signatures}
\label{tab:operator-signatures}
\begin{tabular}{lll}
\toprule
Operator & Operand types & Result type \\
\midrule
\texttt{+}  & Integer $\times$ Integer & Integer \\
\texttt{-}  & Integer $\times$ Integer & Integer \\
\texttt{*}  & Integer $\times$ Integer & Integer \\
\texttt{/}  & Integer $\times$ Integer & Integer \\
\texttt{+.} & Float $\times$ Float     & Float \\
\texttt{-.} & Float $\times$ Float     & Float \\
\texttt{*.} & Float $\times$ Float     & Float \\
\texttt{/.} & Float $\times$ Float     & Float \\
\texttt{==} & Integer $\times$ Integer & Boolean \\
\texttt{==} & Float $\times$ Float     & Boolean \\
\texttt{!=} & Integer $\times$ Integer & Boolean \\
\texttt{!=} & Float $\times$ Float     & Boolean \\
\bottomrule
\end{tabular}
\end{table}

Any operand combination not listed in
Table~\ref{tab:operator-signatures} is rejected during static type checking.
Each operator token maps to exactly one valid operand signature
(or, for \texttt{==} and \texttt{!=}, two same-type signatures), operator
meaning can be determined without overload resolution or implicit type
coercion.

\section{Design Decisions: No Overloading}
\label{sec:no-overloading}
The language definition adopted for this project states: ``No overloading with the usual precedence
of multiplication and division over addition and subtraction. Hence no explicit
declaration.'' This section explains how the language design satisfies this
requirement and why alternatives were rejected.

\textbf{Definition of overloading.}
An operator is overloaded when the same syntactic symbol carries more than one
type signature, requiring the compiler to inspect operand types to determine the
operator's meaning or result type. Under this definition, a language where
\texttt{+} means both integer addition
$(\textit{Integer}, \textit{Integer}) \rightarrow \textit{Integer}$ and
floating-point addition $(\textit{Float}, \textit{Float}) \rightarrow \textit{Float}$
--- or worse, mixed-type addition
$(\textit{Integer}, \textit{Float}) \rightarrow \textit{Float}$ --- overloads the
\texttt{+} symbol.

\textbf{Distinct operator tokens.}
Following the Caml~Light convention, this language assigns a unique token to
each arithmetic meaning. Integer arithmetic uses \texttt{+}, \texttt{-},
\texttt{*}, \texttt{/}; floating-point arithmetic uses \texttt{+.},
\texttt{-.}, \texttt{*.}, \texttt{/.}. Each token has exactly one type
signature. The type of any expression can be determined from the operator token
alone, without inspecting operand types. This is a strict
interpretation of ``no overloading.''

\textbf{Rejected alternative: implicit numeric promotion.}
An alternative design would use a single set of operators and silently promote
\texttt{Integer} operands to \texttt{Float} when mixed with a \texttt{Float}
operand (e.g.\ \texttt{1 + 2.5} $\rightarrow$ \texttt{3.5 : Float}). Proponents
of this approach argue that promotion is not overloading because the operator
retains a single ``arithmetic addition'' meaning. However, this argument is
difficult to sustain formally: the operator \texttt{+} would have three distinct
input-type combinations --- $(\textit{Integer}, \textit{Integer})$,
$(\textit{Integer}, \textit{Float})$, and $(\textit{Float}, \textit{Float})$ ---
with two different result types, a dispatch table that the compiler must resolve
from operand types, which is the hallmark of overloading. The same language
definition also states ``hence no explicit declaration,'' implying that type
inference should be unambiguous; implicit promotion introduces cases where the
inferred result type depends on operand types rather than on the operator itself.

\textbf{Connection to type inference.}
Since every operator token uniquely determines its operand and result types,
the type inference algorithm described in Chapter~5 can determine the type of
any expression without backtracking, global unification, or global constraint
propagation. This
determinism directly satisfies the specification's coupling of ``no
overloading'' with ``no explicit declaration.''

Table~\ref{tab:operator-design-comparison} contrasts the two designs.

\begin{table}[H]
\centering
\caption{Comparison of operator design approaches}
\label{tab:operator-design-comparison}
\small
\begin{tabular}{p{3.2cm}p{4.8cm}p{4.8cm}}
\toprule
\textbf{Expression} & \textbf{This language (separate tokens)} & \textbf{Promotion-based language} \\
\midrule
\texttt{3 + 4}
  & \texttt{7 : Integer} (\texttt{+} $\rightarrow$ Integer)
  & \texttt{7 : Integer} \\
\texttt{3.0 +. 4.0}
  & \texttt{7.0 : Float} (\texttt{+.} $\rightarrow$ Float)
  & N/A (no \texttt{+.} operator) \\
\texttt{3 + 4.0}
  & \texttt{TypeError}: \texttt{+} requires two Integer operands
  & \texttt{7.0 : Float} (implicit promotion) \\
\texttt{10 / 3}
  & \texttt{3 : Integer} (floor division)
  & \texttt{3.333\ldots{} : Float} (or \texttt{3 : Integer}, design-dependent) \\
Type signatures of \texttt{+}
	& One: $(\texttt{Integer},\texttt{Integer}) \rightarrow \texttt{Integer}$
	& Three: $(\texttt{Integer},\texttt{Integer})$, $(\texttt{Integer},\texttt{Float})$, $(\texttt{Float},\texttt{Float})$ \\
Overloaded?
  & No
  & Arguably yes \\
\bottomrule
\end{tabular}
\end{table}
